% ==STYLE-BEGIN==
\documentclass[UTF8,11pt]{ctexart}
\usepackage{amsmath,amssymb,amsthm}
\usepackage[a4paper,margin=2.4cm,top=2.3cm,bottom=2.3cm]{geometry}
\usepackage{fancyhdr}
\usepackage{enumitem}
\usepackage{hyperref}
\hypersetup{hidelinks}

% ---- 版式 ----
\linespread{1.15}
\setlength{\parskip}{0.35em}
\setlength{\parindent}{1.6em}
\setlist{nosep,leftmargin=2em}

% ---- 页眉 ----
\pagestyle{fancy}
\fancyhf{}
\fancyhead[L]{\small\coursename}
\fancyhead[R]{\small\thepage}
\renewcommand{\headrulewidth}{0.4pt}

% ---- 定理环境：中文用正体，标题加粗，按小节共用一条流水号 ----
\newtheoremstyle{cnthm}{0.7em}{0.7em}{\normalfont}{}{\bfseries}{}{0.5em}{}
\theoremstyle{cnthm}
\newtheorem{thm}{定理}[section]
\newtheorem{lem}[thm]{引理}
\newtheorem{prop}[thm]{命题}
\newtheorem{cor}[thm]{推论}
\newtheorem{defn}[thm]{定义}
\newtheorem{rem}[thm]{注}
\newtheorem{exa}[thm]{例}
\newtheorem{conj}[thm]{猜想}
\renewcommand{\proofname}{\bfseries 证明}
\providecommand{\coursename}{}
% ==STYLE-END==
\renewcommand{\coursename}{数学分析\,·\,1028}
\title{数学分析 1028}
\date{}
\begin{document}
\maketitle

\section{反函数的连续性及求导}

\begin{rem}[Recall]
$f:X\to Y$ 是双射，则 $\forall y\in Y,\ \exists!\,x\in X$ s.t. $f(x)=y$。因此可以定义
\[
\varphi:\ Y\to X,\qquad y\longmapsto x ,
\]
称 $\varphi$ 为 $f$ 的反函数。
\end{rem}

\begin{thm}[反函数的存在性、单调性与连续性]
设 $f:[a,b]\to\mathbb{R}$ 连续且严格单调上升，则其反函数 $\varphi:[\alpha,\beta]\to\mathbb{R}$ 也严格单调上升且连续，其中 $\alpha=f(a)$，$\beta=f(b)$。
\end{thm}

\subsection{先说明 $f$ 存在反函数}

从 $f$ 本身的性质入手。$f:[a,b]\to\mathbb{R}$ 连续，可知 $f([a,b])$ 是一个区间（连续函数把区间映成区间），而且应为有界闭区间 $[\alpha,\beta]$。那么
\[
\alpha=\min\{f(x)\mid x\in[a,b]\},\qquad \beta=\max\{f(x)\mid x\in[a,b]\}.
\]
由 $f$ 在 $[a,b]$ 上严格单增可知 $\alpha=f(a)$，$\beta=f(b)$。

此时可以把 $f$ 看成 $[a,b]\to[\alpha,\beta]$ 的映射：函数本身的特性决定 $f$ 一定是满射（每个 $y$ 都有对应的 $x$），而严格单增又说明 $f$ 是单射。$f$ 既是满射又是单射，所以 $f$ 是双射，一定存在反函数
\[
\varphi:\ [\alpha,\beta]\to[a,b],\qquad y\longmapsto x\quad(x\triangleq\varphi(y)),
\]
满足 $f(x)=f\bigl(\varphi(y)\bigr)=y$。

\subsection{$\varphi$ 也严格单调上升}

即要证：$\forall y_1,y_2\in[\alpha,\beta]$，$y_1<y_2\Rightarrow\varphi(y_1)<\varphi(y_2)$。

\begin{proof}
记 $\varphi(y_1)\triangleq x_1$，$\varphi(y_2)\triangleq x_2$，则 $f(x_1)=y_1$，$f(x_2)=y_2$。由于 $f$ 严格单调且 $y_1\neq y_2$，有 $x_1\neq x_2$，于是只能是 $x_1<x_2$ 或 $x_1>x_2$。

若 $x_1>x_2$，据"$f$ 严格单调上升"得 $y_1>y_2$，与题设 $y_1<y_2$ 矛盾，舍去。

故 $x_1<x_2$，即 $\varphi(y_1)<\varphi(y_2)$。
\end{proof}

\subsection{$\varphi$ 在 $(\alpha,\beta)$ 上连续}

（若要证在整个 $[\alpha,\beta]$ 上连续，只需再证 $\varphi$ 在 $\alpha$ 处右连续、在 $\beta$ 处左连续，做法完全类似。）

\begin{proof}
设 $y_0\in(\alpha,\beta)$，记 $x_0\triangleq\varphi(y_0)\in(a,b)$。任给 $\varepsilon>0$。

\textbf{第一步：把 $\varepsilon$ 缩小到"不越界"。} 因为 $x_0$ 是开区间 $(a,b)$ 的内点，令
\[
\varepsilon_1=\min\left\{\varepsilon,\ \frac{x_0-a}{2},\ \frac{b-x_0}{2}\right\}>0 ,
\]
则 $x_0\pm\varepsilon_1\in(a,b)$，并且 $\varepsilon_1\leq\varepsilon$。

\textbf{第二步：由 $\varepsilon_1$ 造两个点。} 令
\[
x_1=x_0-\varepsilon_1,\qquad x_2=x_0+\varepsilon_1 ,
\]
则 $a\leq x_1<x_0<x_2\leq b$。由 $f$ 严格单增，
\[
f(x_1)<f(x_0)=y_0<f(x_2).
\]

\textbf{第三步：把 $\delta$ 定出来。} 令
\[
\delta=\min\bigl\{f(x_2)-y_0,\ y_0-f(x_1)\bigr\}>0 .
\]
若 $|y-y_0|<\delta$，则
\[
f(x_1)=y_0-\bigl(y_0-f(x_1)\bigr)\leq y_0-\delta<y<y_0+\delta\leq y_0+\bigl(f(x_2)-y_0\bigr)=f(x_2),
\]
即 $f(x_1)<y<f(x_2)$。这时 $y\in\bigl(f(x_1),f(x_2)\bigr)\subseteq(\alpha,\beta)$，所以 $\varphi(y)$ 有定义。

\textbf{第四步：用 $\varphi$ 单增回代。} 对 $f(x_1)<y<f(x_2)$ 两边作用严格单增的 $\varphi$，得
\[
x_1<\varphi(y)<x_2 ,
\]
即
\[
\bigl|\varphi(y)-\varphi(y_0)\bigr|=\bigl|\varphi(y)-x_0\bigr|<\varepsilon_1\leq\varepsilon .
\]
故 $\varphi$ 在 $y_0$ 处连续。
\end{proof}

\begin{rem}[关于"取 $\varepsilon$ 足够小"]
上面的第一步不能省。连续性要求对\textbf{任意} $\varepsilon>0$（包括大到使 $x_0+\varepsilon>b$ 的 $\varepsilon$）都能找到 $\delta$；而第二、三步要求 $x_0\pm\varepsilon\in(a,b)$。把 $\varepsilon$ 换成 $\varepsilon_1\leq\varepsilon$ 就同时满足两边：小 $\varepsilon$ 成功了，大 $\varepsilon$ 自然更宽松。这一步常被一句话"取 $\varepsilon$ 足够小"带过，严格写应如上处理。
\end{rem}

\begin{rem}[依赖方向，不是循环]
$\varepsilon\longrightarrow(x_1,x_2)\longrightarrow\delta$ 是单向的：$\varepsilon$ 是外面给的，$x_1,x_2$ 由 $\varepsilon$ 定，$\delta$ 再由 $x_1,x_2$ 定，$\delta$ 允许依赖 $\varepsilon$。这正是 $\varepsilon$-$\delta$ 语言所要求的。
\end{rem}

\section{反函数的导数}

\begin{thm}[反函数求导]
设 $f$ 在 $x_0$ 附近严格单增且连续，又 $f'(x_0)\neq0$，则其反函数 $x=\varphi(y)$ 在 $y_0=f(x_0)$ 处可导，且
\[
\varphi'(y_0)=\frac{1}{f'(x_0)} .
\]
\end{thm}

\begin{rem}[为什么"局部"就够]
不必要求 $f$ 在整个定义域上单调，只要在 $x_0$ 附近严格单增且连续就够了。具体说，取 $\eta>0$ 使 $(x_0-\eta,x_0+\eta)\subseteq(a,b)$，则 $f$ 在其上严格单增且连续；由介值定理，$f\bigl((x_0-\eta,x_0+\eta)\bigr)$ 是一个区间，又因严格单增，端点取不到，所以它是\textbf{开区间}
\[
\bigl(f(x_0-\eta),\ f(x_0+\eta)\bigr)\ni y_0 .
\]
在这段开区间上 $f$ 严格单增且是满射，因而是双射，$\varphi$ 在 $y_0$ 的某个邻域内有定义、连续、严格单增。所以"$\varphi$ 在 $y_0$ 处可导"是一个\textbf{局部}结论：无论 $f$ 在远处怎样，只要 $f'(x_0)\neq0$（从而 $f$ 在 $x_0$ 附近严格单调），总能在 $y_0$ 附近找到一段有反函数的区间。
\end{rem}

\begin{proof}
\[
\varphi'(y_0)=\lim_{y\to y_0}\frac{\varphi(y)-\varphi(y_0)}{y-y_0}
\overset{y=f(x)}{=}
\lim_{x\to x_0}\frac{\varphi\bigl(f(x)\bigr)-\varphi\bigl(f(x_0)\bigr)}{f(x)-f(x_0)}
=\lim_{x\to x_0}\frac{x-x_0}{f(x)-f(x_0)} ,
\]
其中换元 $y=f(x)$ 是合法的：$f$ 与 $\varphi$ 都在相应点连续且互为反函数，故 $y\to y_0\iff x\to x_0$。

又
\[
\frac{x-x_0}{f(x)-f(x_0)}=\frac{1}{\dfrac{f(x)-f(x_0)}{x-x_0}} ,
\]
由于 $f'(x_0)\neq0$，可以用四则运算（商的极限），得
\[
\varphi'(y_0)=\frac{1}{\lim\limits_{x\to x_0}\dfrac{f(x)-f(x_0)}{x-x_0}}=\frac{1}{f'(x_0)} . \qedhere
\]
\end{proof}

\begin{cor}
在定理 2.1 的条件下，若把 $y=f(x)$、$x=\varphi(y)$ 写回 $y$ 的函数，则
\[
\varphi'(y)=\frac{1}{f'\bigl(\varphi(y)\bigr)} ,
\]
即反函数的导数等于原函数导数在对应点的倒数。
\end{cor}

\end{document}
